% Generated draft: requires mathematical and source review.
\documentclass[11pt]{article}
\usepackage[margin=25mm]{geometry}
\usepackage{fontspec}
\setmainfont{Latin Modern Roman}
\usepackage{amsmath,amssymb,mathtools,unicode-math}
\setmathfont{Latin Modern Math}
\usepackage{enumitem,needspace,xurl,hyperref}
\setlength{\parindent}{0pt}
\setlength{\parskip}{5pt}
\setlength{\emergencystretch}{3em}
\newcommand{\N}{\mathbb{N}}
\newcommand{\Z}{\mathbb{Z}}
\newcommand{\Q}{\mathbb{Q}}
\newcommand{\R}{\mathbb{R}}
\newcommand{\C}{\mathbb{C}}
\newcommand{\F}{\mathbb{F}}
\newcommand{\A}{\mathbb{A}}
\newcommand{\PP}{\mathbb{P}}
\newcommand{\E}{\mathbb{E}}
\newcommand{\eps}{\varepsilon}
\newcommand{\dd}{\,\mathrm{d}}
\begin{document}
\section*{UnsolvedMath ID: 701594}
\textbf{Source problem code:} catalogue entry
\subsection{Definitions and mathematical statement}
Let \(G\) denote the group written \(\mathrm{CT}(\mathbb{Z})\) in the supplied source. The source refers to Kourovka Notebook Problem 17.57 for its definition, but that definition is absent from the supplied record. Thus the concrete meaning of \(G\) remains unresolved here.

A group \(H\) is periodic if every element has finite order, that is, if \(\forall h\in H\;\exists n\in\mathbb{Z}_{\geq 1}\; h^n=e_H\). It is finitely generated if \(H=\langle h_1,\ldots,h_r\rangle\) for some positive integer \(r\) and elements \(h_1,\ldots,h_r\in H\).

The question attributed to S. Kohl (2018) is: does there exist a subgroup \(H\leq G\) such that \(H\) is infinite, finitely generated, and periodic? Equivalently, is there an infinite subgroup \(H\leq G\) satisfying
\[
\bigl(\exists r\in\mathbb{Z}_{\geq 1}\;\exists h_1,\ldots,h_r\in H:\ H=\langle h_1,\ldots,h_r\rangle\bigr)
\quad\text{and}\quad
\bigl(\forall h\in H\;\exists n\in\mathbb{Z}_{\geq 1}:\ h^n=e_H\bigr)?
\]
The exponent \(n\) may depend on \(h\); a common finite exponent is not required. This formulation requires the missing definition from Problem 17.57 to identify the ambient group.
\subsection{Short English statement}
Does the source-designated group \(\mathrm{CT}(\mathbb{Z})\) contain an infinite subgroup generated by finitely many elements in which every element has finite order? The supplied record omits the referenced definition of the ambient group.
\subsection{Sources}
\begin{itemize}
\item Supplied local UnsolvedMath registry record, permanent ID \(701594\), exact source code \texttt{catalogue entry}, titled \emph{Kourovka Notebook Problem 19.46}. Its status is supplied as \texttt{open}; this is not an independently verified assessment.
\item \emph{The Kourovka Notebook (Unsolved Problems in Group Theory)}, 19th issue (2018), Problem 19.46; also identified in the supplied record with the 21st issue edition (2026), arXiv:\texttt{1401.0300v48}. Proposed by S. Kohl in 2018. The supplied collection description names E. I. Khukhro and V. D. Mazurov as editors of Issues 1--21 (1965--2026).
\item Kourovka Notebook Problem 17.57 is explicitly referenced by the statement for the definition of \(\mathrm{CT}(\mathbb{Z})\). Its contents are not supplied and have not been independently consulted.
\end{itemize}
\end{document}
